\documentclass[12pt]{article}%
\usepackage{amsfonts}
\usepackage{fancyhdr}
\usepackage{comment}
\usepackage[a4paper, top=2.5cm, bottom=2.5cm, left=2.2cm, right=2.2cm]%
{geometry}
\usepackage{times}
\usepackage{amsmath}
\usepackage{changepage}
\usepackage{amssymb}
\usepackage{graphicx}

\setcounter{MaxMatrixCols}{30}
\newtheorem{theorem}{Theorem}
\newtheorem{acknowledgement}[theorem]{Acknowledgement}
\newtheorem{algorithm}[theorem]{Algorithm}
\newtheorem{axiom}{Axiom}
\newtheorem{case}[theorem]{Case}
\newtheorem{claim}[theorem]{Claim}
\newtheorem{conclusion}[theorem]{Conclusion}
\newtheorem{condition}[theorem]{Condition}
\newtheorem{conjecture}[theorem]{Conjecture}
\newtheorem{corollary}[theorem]{Corollary}
\newtheorem{criterion}[theorem]{Criterion}
\newtheorem{definition}[theorem]{Definition}
\newtheorem{example}[theorem]{Example}
\newtheorem{exercise}[theorem]{Exercise}
\newtheorem{lemma}[theorem]{Lemma}
\newtheorem{notation}[theorem]{Notation}
\newtheorem{problem}[theorem]{Problem}
\newtheorem{proposition}[theorem]{Proposition}
\newtheorem{remark}[theorem]{Remark}
\newtheorem{solution}[theorem]{Solution}
\newtheorem{summary}[theorem]{Summary}
\newenvironment{proof}[1][Proof]{\textbf{#1.} }{\ \rule{0.5em}{0.5em}}

\newcommand{\Q}{\mathbb{Q}}
\newcommand{\R}{\mathbb{R}}
\newcommand{\C}{\mathbb{C}}
\newcommand{\Z}{\mathbb{Z}}

%%%%
% ME %
%%%%

\usepackage{listings}
\usepackage{color}
\usepackage{xcolor}
\usepackage{mdframed}

\definecolor{dkgreen}{rgb}{0,0.6,0}
\definecolor{gray}{rgb}{0.5,0.5,0.5}
\definecolor{mauve}{rgb}{0.58,0,0.82}

\lstset{%frame=tb,
language=Matlab,
aboveskip=3mm,
belowskip=3mm,
showstringspaces=false,
columns=flexible,
basicstyle={\small\ttfamily},
numbers=none,
numberstyle=\tiny\color{gray},
keywordstyle=\color{blue},
commentstyle=\color{dkgreen},
stringstyle=\color{mauve},
breaklines=true,
breakatwhitespace=true,
tabsize=3
}

\usepackage{outlines}
\usepackage{enumitem}
%\setenumerate[1]{label=\Roman*.}
%\setenumerate[2]{label=\Alph*.}
%\setenumerate[3]{label=\roman*.}
%\setenumerate[4]{label=\alph*.}

\usepackage{titlesec} 

\titleformat{\subsection}[runin]
  {\normalfont\large\bfseries}{\thesubsection}{1em}{}
\titleformat{\subsubsection}[runin]
  {\normalfont\normalsize\bfseries}{\thesubsubsection}{1em}{}

% for enumerate spacing
\newenvironment{tight_enum}{
\begin{enumerate}[label=\alph*.]
\setlength{\itemsep}{0pt}
\setlength{\parskip}{0pt}
}{\end{enumerate}}

\usepackage{exercise}

\setlength\parindent{0pt}

\newcommand{\code}[1]{\texttt{#1}}
\newcommand{\shp}{$>>>$ }
\newcommand{\tab}{\hspace*{45pt}}

\usepackage{multicol}
\setlength{\columnsep}{0.6cm}

\begin{document}

\title{Homework 09: \\ NumPY and Matplotlib}

\author{EEC 3150}
\date{}

\maketitle

\begin{center}
\fbox{\fbox{\parbox{5.5in}{\centering
For each Python exercise, write a separate Python script. Submit all scripts on Blackboard under the appropriate assignment. Name your script files with the following format: \\
\vspace{10pt}
HW$<$assignment number$>$\_Ex$<$exercise number$>$\_$<$FirstInitial$><$LastName$>$.py \\
\vspace{10pt}
Example: HW01\_Ex01\_GWest.py \\
\vspace{10pt}
For short answer or math exercises, either do your work on a sheet of paper and upload a picture of it to Blackboard or submit some kind of text file or Word document to Blackboard. You can place all such exercises in a single document as long as you clearly label each exercise. These files should follow a similar naming convention as the Python scripts.
}}}
\end{center}


\pagebreak


\begin{multicols*}{2}



%%% NEW EXERCISE %%%
\begin{Exercise}[origin={10 points, Python}]

\begin{enumerate}[leftmargin=*]
	\item Write a function \code{Boundary(n,value)} which will create an $n\times n$ matrix in which the elements on the top, bottom, left, and right boundaries \code{value} while the inside values are equal to 0.
	\item Demonstrate your function by calling it and printing the output
\end{enumerate}

Here is an example: \code{Boundary(4,1)} should produce: \\

\code{[[1 1 1 1] \\[0pt]
[1 0 0 1] \\[0pt]
[1 0 0 1] \\[0pt]
[1 1 1 1]]
}

\end{Exercise}

%%% NEW EXERCISE %%%
\begin{Exercise}[origin={10 points, Python}]

\begin{enumerate}[leftmargin=*]
	\item Write a function \code{Truncate(arr)} which an array of any shape and return a new array which is a copy of the old array with the following modifications:
	\begin{itemize}
		\item Set all values $<0$ to be $=0$
		\item Set all values $>0$ to be $=1$
	\end{itemize}
	\item Demonstrate your function by calling it and printing the output
\end{enumerate}

Here is an example: \\ \code{Truncate(np.array([-1,2,10,-4]))} should produce: \\

\code{[0 1 1 0]}

\end{Exercise}

%%% NEW EXERCISE %%%
\begin{Exercise}[origin={20 points, Python}]

The exercise will have you generate random data, plot it, and index it.

\begin{enumerate}[leftmargin=*]
	\item Generate random data
	\begin{itemize}[leftmargin=*]
		\item Create 1D \code{x,y} arrays of 1000 random numbers between -2 and 2
	\end{itemize}
	\item Plot the data
	\begin{itemize}[leftmargin=*]
		\item Plot the \code{x,y} arrays as blue dots
	\end{itemize}
	\item Index the data
	\begin{itemize}[leftmargin=*]
		\item Through whatever means you choose, create \code{x2,y2} arrays which only havve the points where $-1.5<x<1$ and $-1<y<0.9$
	\end{itemize}
	\item Plot the new data
	\begin{itemize}[leftmargin=*]
		\item Plot the \code{x2,y2} arrays on the same plot as red x's with an increased \code{markersize} for readability
	\end{itemize}
\end{enumerate}

\end{Exercise}

\vfill\null
\columnbreak



\vfill\null
\end{multicols*}




\end{document}






